% 引言章节
% 说明：
% 1. 本文件设计为被主文档 \input{} 进来使用。
% 2. 内容来源于 04_paper.md 的 Introduction 部分。
% 3. 这里只做 Markdown 到 LaTeX 的语法迁移，不改写正文内容。

\section{Introduction}

基于辐射场的同步定位与建图（SLAM）系统近年来受到广泛关注。其中，3D Gaussian Splatting（3DGS）以其显式高斯原语表示和高效的可微光栅化渲染，在渲染速度和优化效率上展现出显著优势，已成为在线照片级场景重建的重要基础表示。在此基础上，以 Gaussian-LIC 系列为代表的 LiDAR-Inertial-Camera 融合系统进一步证明了多传感器条件下增量式高斯地图在鲁棒性和渲染质量上的潜力。

然而，现有增量式高斯 SLAM 系统在地图构建质量上仍面临若干根本性的局限。

第一，新高斯的初始化依赖手工规则。当新关键帧到达时，系统通常基于 LiDAR 点或深度补全结果生成候选点，然后以固定规则赋予其初始属性——尺度由深度线性推算、不透明度设为固定低值、颜色取单帧 DC 分量、旋转无约束或仅对齐视线方向。这种手工初始化往往质量有限，导致新高斯需要经历较长的在线优化才能收敛到合理状态，在快速运动或低纹理区域尤为脆弱。

第二，体高斯（3DGS）表示缺乏几何约束。标准 3DGS 中的高斯原语为三维椭球体，不具有明确的表面法线和平面几何意义。这限制了深度渲染的几何一致性，也使得几何先验（如单目法线估计）难以与高斯表示有效耦合。2D Gaussian Splatting（2DGS）通过将高斯压缩为嵌入三维空间的二维有向盘面（surfel），赋予每个原语明确的法线方向和切平面结构，在表面贴合和几何质量上已取得显著进步，但尚未被系统性地引入增量式 SLAM 场景。

第三，现有学习型先验的利用方式单一。尽管 Gaussian-LIC2 等系统已引入深度补全网络以缓解 LiDAR 覆盖不足问题，但学习型先验仅被用于确定补点位置，尚未扩展到高斯属性（尺度、方向、颜色、不透明度）的全面初始化。与此同时，pixelSplat、MVSplat 等前馈高斯预测方法已在离线场景中证明了从图像特征直接回归高斯属性的可行性，但这些方法面向一次性多视图重建，而非增量式 SLAM 中的逐帧补点场景。

针对上述问题，本文提出了一种面向增量式 2DGS 场景重建的上下文感知前馈高斯回归系统。系统在外部稳定前端提供位姿的条件下，以 2DGS surfel 为地图基础表示，通过融合 SPNet 深度补全与 DA3（Depth Anything）法线先验确定候选点及其几何方向，进而由一个前馈回归网络直接预测新高斯的完整属性，随后新高斯纳入在线光度优化迭代。该系统还包含一套完整的自蒸馏训练闭环，使前馈网络能在无需外部标注数据集的条件下训练和部署。

本文的主要贡献如下：

\begin{itemize}
    \item 基于 2DGS surfel 的几何感知增量式高斯地图后端。不同于现有以 3DGS 体高斯为默认表示的增量式系统，本文以 2DGS 的表面化表示为基础，围绕 surfel 的法线语义构建了融合 SPNet 深度补全与 DA3 法线先验的统一几何管线，并通过深度畸变与法线一致性约束驱动地图更新，为后续的前馈高斯插入方法提供了几何可解释的表示基础。

    \item 上下文感知的前馈高斯插入方法。本文提出一种用于替代传统手工初始化的前馈高斯属性回归方法。该方法对每个候选点联合利用当前帧与历史帧的原始与渲染图像特征，通过逐点跨注意力融合多视角上下文，并以 DA3 法线先验为锚点、经 Gram-Schmidt 正交化构造 surfel 旋转，在单次前向传播中直接回归完整的 2DGS surfel 属性。

    \item 面向在线增量建图的自蒸馏训练闭环。本文构建了一套使前馈插入方法在无需外部标注数据集条件下完成训练的自蒸馏管线。该闭环以高斯为核心单位构造训练样本，通过因果正确的回溯特征采集避免信息泄露，并显式消解 2DGS 固有的尺度-旋转歧义，实现从在线建图、离线训练到重新部署的自增强过程。
\end{itemize}
